\begin{FSChapterPage}{6}{Funcionamento de Equipa}
\label{toc:6}
\FSSectionStart[toc:6.1]{6.1 Sanções}
\FSRuleA{6.1.1}{Podem ser aplicadas sanções a qualquer membro da equipa, cuja interpretação fica ao cargo do mesmo e Direção.}
\FSRuleA{6.1.2}{O não cumprimento com o Regulamento Geral ou demais legislação aplicável, pode levar à atribuição de Sanção.}
\FSRuleA{6.1.3}{Pode ser apresentado um protesto à sanção e suas consequências num período de tempo de até uma semana após conhecimento da mesma.}
\FSRuleA{6.1.4}{Uma sanção constitui um aviso, apontando o incumprimento supramencionado ao membro em causa.}

\FSSection[toc:6.2]{6.2 Demissões}
\FSRuleA{6.2.1}{Mediante sanções e votação favorável da Direção segundo a Metodologia de Decisão, um membro da equipa poderá ser demitido com efeito imediato.}
\FSRuleA{6.2.2}{A demissão de qualquer membro individual da Direção poderá implicar a cooptacão por parte da própria Direção de um membro da equipa.}
\FSRuleA{6.2.3}{Um descontentamento demonstrado pela equipa poderá levar à destituição da Direção, mediante deliberação de pelo menos dois terços dos membros da equipa a aprovar.}
\FSRuleA{6.2.4}{A demissão do Team Leader implicará um novo processo eleitoral, que determinará a Direção sucessora. Este processo eleitorial deverá ocorrer num prazo máximo de quinze dias.}

\FSSection[toc:6.3]{6.3 Incompatibilidades}
\FSRuleA{6.3.1}{A presença de um Member noutro Núcleo ou Associação deve ser revista por parte da Direção de equipa, decidindo esta a metodologia a seguir.}
\end{FSChapterPage}

\begin{FSBodyPage}
\FSRuleA{6.3.2}{Qualquer situação de potencial incompatibilidade é revista e decidida pela Direção.}
\FSRuleA{6.3.3}{É incompatível a presença noutro Núcleo ou Associação, de um membro da Direção ou Board da equipa.}

\FSSection[toc:6.4]{6.4 Recrutamento}
\FSRuleA{6.4.1}{A equipa segue um processo de recrutamento público, aberto aos estudantes de primeiro e segundo ciclos da Universidade do Porto.}
\FSRuleA{6.4.2}{A definição e anúncio do calendário de recrutamento é da responsabilidade da Board da equipa.}
\FSRuleA{6.4.3}{O recrutamento visa assegurar que a continuidade e manutenção de equipa são asseguradas, procurando que esta cumpra com os seus objetivos e Visão.}
\FSRuleA{6.4.4}{A qualidade e transparência do processo de recrutamento devem ser procurados, seguindo sempre critérios claros e rigorosos.}

\FSSection[toc:6.5]{6.5 Instalações}
\FSRuleA{6.5.1}{As instalações utilizadas pela equipa devem ser cuidadosamente utilizadas, procurando a sua conservação:}
\FSRuleB{6.5.1.1}{Os danos causados intencionalmente ou de forma irresponsável, às instalações serão reparados/pagos pelos responsáveis pelo ato.}
\FSRuleB{6.5.1.2}{Não é permitido o uso das instalações para atos que não sejam relacionados com a sua finalidade.}
\FSRuleB{6.5.1.3}{Os membros, enquanto utilizadores do espaço, devem cumprir as regras de funcionamento estabelecidas:}
\FSRuleC{6.5.1.3.1}{Respeitar a integridade do espaço de trabalho.}
\FSRuleC{6.5.1.3.2}{Respeitar os demais membros a trabalhar.}
\FSRuleC{6.5.1.3.3}{Procurar não perturbar possíveis reuniões que existam em simultâneo no espaço.}
\FSRuleB{6.5.1.4}{Qualquer violação do Regulamento Geral resulta no afastamento do responsável ou responsáveis pela mesma e respetiva avaliação da situação pela Direção.}
\FSRuleA{6.5.2}{A equipa dispõe de um local de trabalho, a sala B320.}
\FSRuleA{6.5.3}{É possível de incluir material pessoal nas instalações utilizadas pela equipa, desde que:}
\FSRuleB{6.5.3.1}{Este não cause incómodo no demais membros da equipa.}
\FSRuleB{6.5.3.2}{O espaço ocupado por este seja adequado.}
\FSRuleB{6.5.3.3}{Houver possibilidade de existência de incompatibilidade com algum material pré-existente na instalação.}
\end{FSBodyPage}

\begin{FSBodyPage}
\FSSectionStart[toc:6.6]{6.6 Metodologia de decisão}
\FSRuleA{6.6.1}{A metodologia de decisão seguida pela Direção e Board, abaixo explicada, carece de um documento escrito sempre que utilizada, contendo a decisão a ser tomada assim como assinatura dos membros intervenientes sob o resultado da votação.}
\FSRuleA{6.6.2}{Cada voto conta como um ponto, tendo qualquer membro da Direção ou Board o mesmo peso nesse voto:}
\FSRuleB{6.6.2.1}{Um voto favorável conta como um ponto positivo (+1).}
\FSRuleB{6.6.2.2}{Uma abstenção conta como um ponto nulo (=0).}
\FSRuleB{6.6.2.3}{Um voto contra conta como um ponto negativo (-1).}
\FSRuleA{6.6.3}{O somatório da votação dita o resultado da mesma:}
\FSRuleB{6.6.3.1}{Se for superior a zero, a decisão é favorável (\textit{>} 0).}
\FSRuleB{6.6.3.2}{Se for igual a zero e a abstenção for inferior a um terço dos membros, considera-se um empate (=0).}
\FSRuleB{6.6.3.3}{Se for inferior a zero, a decisão é desfavorável (\textit{<} 0).}
\FSRuleA{6.6.4}{Na eventualidade de um empate, o Team Leader tem direito a voto de qualidade.}
\FSRuleA{6.6.5}{É necessário voto da totalidade dos membros da Direção para a votação ser contabilizada como válida, exceto se houver impossibilidade prolongada de exercer esse dever.}
\FSRuleA{6.6.6}{Certas decisões, se considerado necessário, precisam de somatório favorável da Direção, Board e eventualmente Team Leader, para serem aprovadas.}
\FSRuleA{6.6.7}{A tipologia de decisões no anterior ponto mencionadas cabe à Board de definir.}
\FSRuleA{6.6.8}{A confidencialidade da votação e seu resultado, se necessário, deve ser assegurada.}

\FSSection[toc:6.7]{6.7 Dissolução}
\FSRuleA{6.7.1}{A equipa pode ser dissolvida, se assim decidido pela Direção da mesma, devido a:}
\FSRuleB{6.7.1.1}{Impossibilidade financeira de manter a equipa.}
\FSRuleB{6.7.1.2}{Inexistência de membros.}
\FSRuleB{6.7.1.3}{Inexistência prolongada de atividade.}
\FSRuleB{6.7.1.4}{Manifestado interesse na sua execução com base em baixo assinado por 75\% dos membros da equipa.}
\FSRuleA{6.7.2}{Na eventualidade de uma dissolução, o património da equipa ficará sob tutela da Comunidade.}
\end{FSBodyPage}

\begin{FSBodyPage}
\FSSectionStart[toc:6.8]{6.8 Autonomização de Entidade}
\FSRuleA{6.8.1}{Por decisão da Board da equipa, pode ser deliberada a autonomização da FS FEUP como entidade jurídica distinta da AEFEUP, processo que terá início mediante a apresentação formal desta intenção via comunicação escrita à Direção da AEFEUP.}
\FSRuleA{6.8.2}{Após esta comunicação, será realizada a autonomização da entidade FS FEUP, que culminará na criação de um nova associação, dando seguimento aos projetos em curso e cumprindo todos os formalismos legais.}
\FSRuleA{6.8.3}{Uma vez constituída a associação, a AEFEUP deverá transferir a totalidade do património que detenha da equipa, num prazo máximo de quinze dias úteis.}

\FSSection[toc:6.9]{6.9 Receitas e Despesas}
\FSRuleA{6.9.1}{A equipa tem um orçamento próprio, gerido pela Board da mesma.}
\FSRuleA{6.9.2}{As receitas da equipa advêm de:}
\FSRuleB{6.9.2.1}{Atribuição de subsídios por parte dos parceiros institucionais da equipa.}
\FSRuleB{6.9.2.2}{De heranças, legados e outros subsídios a estes atribuídos.}
\FSRuleB{6.9.2.3}{De patrocínios e apoios de qualquer entidade ou instituição.}
\FSRuleB{6.9.2.4}{De liberalidades aceites pela equipa.}
\FSRuleB{6.9.2.5}{Da prestação de bens, serviços ou demais iniciativas.}
\FSRuleA{6.9.3}{As despesas da equipa devem ser geridas pela Board e Financial Manager, sendo utilizadas única e exclusivamente para colmatar as necessidades da equipa.}

\FSSection[toc:6.10]{6.10 Conduzir o protótipo}
\FSRuleA{6.10.1}{Carece de uma aprovação por parte da Direção de equipa.}
\FSRuleA{6.10.2}{Requer a posse de Carta de Condução do Tipo B.}
\FSRuleA{6.10.3}{Requer a capacidade de manusear o protótipo em segurança, sendo necessária uma formação e treino promovida pela equipa.}
\FSRuleA{6.10.4}{É necessário pertencer à equipa para conduzir o protótipo em competição.}
\FSRuleA{6.10.5}{Qualquer custo associado a um dano ao protótipo, seja intencional, ou devido a irresponsabilidade, promovido pelo piloto, deve ser acarretado por este.}
\end{FSBodyPage}

\begin{FSBodyPage}
\FSSectionStart[toc:6.11]{6.11 Participação em competições}
\FSRuleA{6.11.1}{A curadoria do processo de candidatura a competições é da responsabilidade da Board.}
\FSRuleA{6.11.2}{O processo de decisão acerca de participação em competições deve ter em conta:}
\FSRuleB{6.11.2.1}{Disponibilidade do membro.}
\FSRuleB{6.11.2.2}{Interesse do membro.}
\FSRuleB{6.11.2.3}{Interesse global da equipa: continuidade.}
\FSRuleB{6.11.2.4}{Orçamento da equipa.}
\FSRuleA{6.11.3}{A participação em competições pode requerer a comparticipação no orçamento da equipa, de apoio dos membros.}
\FSRuleA{6.11.4}{Uma vez divulgada a lista de participantes nas competições, cada membro tem até uma semana para apresentar um protesto, se necessário.}
\FSRuleA{6.11.5}{O apoio monetário não será reembolsado, independentemente do motivo que causou a necessidade, uma vez feitas as compras que concernem a participação do membro em competições.}
\FSRuleA{6.11.6}{É requerido no período de participação que cada membro cumpra com o regulamento da própria competição, Regulamento Geral, Protocolos e demais legislação aplicável.}

\FSSection[toc:6.12]{6.12 Manuseamento de património}
\FSRuleA{6.12.1}{O património da equipa pode ser utilizado para desenvolvimento do projeto, desde que aprovado pela Direção.}
\FSRuleA{6.12.2}{A utilização carece da escrita de um documento prévio, contendo uma descrição, objetivos e lugar para posterior anotação dos resultados.}
\FSRuleA{6.12.3}{A aquisição de material a ser utilizado para o desenvolvimento feito pela equipa é passível de ser feita, desde que justificado e enquadrado no orçamento global.}
\FSRuleA{6.12.4}{Qualquer custo associado a um dano intencional ao património, promovido pelo membro que o usou deve ser acarretado por este.}
\FSRuleA{6.12.5}{Um componente da posse de um membro, uma vez implementado no protótipo, carece de uma autorização da Direção para ser removido.}
\FSRuleA{6.12.6}{Os componentes resultantes de um projeto em ambiente de Dissertação ou Projeto Integrador/Projeto Experimental \& Computacional, quando elaborados em colaboração com a equipa, pertencentem à mesma.}
\end{FSBodyPage}

\begin{FSBodyPage}
\FSSectionStart[toc:6.13]{6.13 Épocas Especiais}
\FSRuleA{6.13.1}{Pode ser proposto à Direção da FEUP o acesso de um Membro a Épocas Especiais.}
\FSRuleA{6.13.2}{A quantidade de acessos por Época é definida pelo Faculty Advisor e Team Leader.}
\FSRuleA{6.13.3}{O Membro interessado deverá informar o seu Department Leader da necessidade de proposta.}
\FSRuleA{6.13.4}{Carece de uma aprovação por parte da Direção de equipa, com base em:}
\FSRuleB{6.13.4.1}{Proatividade e impacto do Membro individual.}
\FSRuleB{6.13.4.2}{Análise macro da comparação dos interessados em relação às vagas disponíveis.}
\end{FSBodyPage}
